\begin{abstract}
\TODO{Translate and condense the existing German abstract (\textasciitilde300 words) to
\textasciitilde200 words. Final numbers and findings can only be added after the
experiments; until then, the abstract describes goal and method only.}

Large Language Models (LLMs) are increasingly used to generate source code, yet their
evaluation focuses predominantly on functional correctness. The energy consumed when
the generated code is \emph{executed} has received comparatively little attention.
This paper presents an empirical study of the energy efficiency of LLM-generated code
along three factors: the generating LLM (Claude, GPT, and Gemini, accessed via their
official coding CLIs), the target programming language (Python, Java,
C++), and the hardware platform on which the code runs (laptop, Raspberry~Pi,
server). Using a full factorial design over \TBD{} programming tasks, we measure energy
consumption, execution time, and peak memory usage of functionally correct solutions
and analyse main and interaction effects with mixed-effects models. In particular, we
investigate whether the relative ranking of LLMs with respect to energy efficiency is
stable across languages and platforms.
\TODO{Results and implications.}
\end{abstract}

\begin{IEEEkeywords}
Large language models, code generation, energy efficiency, green software engineering,
empirical study, programming languages, Raspberry~Pi
\end{IEEEkeywords}
